\chapter{多普勒效应}\label{ch:g12:doppler-effect}

一辆救护车驶近，警笛大作，呼啸而过——就在它经过的那一瞬间，
音调明显地向下滑落，尽管司机什么也没碰。警笛本身是诚实的，波却不然：
运动把它在前方挤压，在后方拉伸。本章将度量这种挤压，
再把同一思想从警笛推广到雷达测速枪、动脉中的血流，
乃至一颗因星光中微小晃动而暴露行踪的行星。

\section{效应}

\begin{definition}[多普勒效应]\label{def:g12:doppler-effect:doppler}
\emph{多普勒效应}\index{多普勒效应}是指波的接收频率
在波源与接收器相对运动时发生的变化：
两者靠近时接收频率 $f'$ 高于发出频率 $f$，远离时则更低，
仅当两者距离在该瞬间不再变化时，接收频率恰为 $f$。
\end{definition}

\begin{remark}[音调而非响度]\label{rem:g12:doppler-effect:pitch}
救护车经过时，有两件事在变化：\emph{响度}的增强与减弱
是因为波随距离扩散（\cref{ch:g12:sound-acoustics}）；\emph{音调}的变化
则源于靠近或远离的速率，并在经过的瞬间突然下降。
响度说明“多远”；音调说明“多快”。
\end{remark}

\section{运动的波源：拥挤的波前}

关键图像是：波峰一旦离开波源，便由介质“接管”。每个波峰以传播速度 $v$
（\cref{ch:g12:mechanical-waves}）扩展为圆，圆心位于它被
\emph{发出}时所在的位置——介质既不知道，也不关心波源已经向前移动。

\begin{proposition}[运动的波源]\label{prop:g12:doppler-effect:moving-source}
波源以周期 $T$（频率 $f = 1/T$）发出波峰，
并以速率 $v_s < v$ 沿直线运动。介质中波的传播速度为 $v$。
介质中静止的观察者接收到
\[
\lambda' = (v \mp v_s)\,T, \qquad f' = f\,\frac{v}{v \mp v_s},
\]
上面的符号对应靠近的波源，下面的符号对应远离的波源。
\end{proposition}

\begin{proof}
两相邻波峰之间，第一个波峰前进 $vT$，而波源前进 $v_s T$。前方的
第二个波峰的发出点因此离第一个波峰 $v_s T$ \emph{更近}，即 $\lambda' = (v - v_s)T$；
后方则增加同样的 $v_s T$。随后，波峰掠过静止的观察者，
其速率为 $v$，每隔 $\lambda'/v$ 秒经过一个，故 $f' = v/\lambda' = f v/(v \mp v_s)$。
\end{proof}

\begin{omfigure}
\begin{tikzpicture}[font=\small]
  \foreach \i/\r in {0/3.0, 1/2.25, 2/1.5, 3/0.75} {
    \draw[omDef] ({0.3*\i},0) circle ({\r});
    \fill[gray] ({0.3*\i},0) circle (1.1pt);
  }
  \fill (1.2,0) circle (2pt);
  \draw[-Stealth, omThm, thick] (1.2,0) -- (1.75,0)
    node[below right=-2pt] {$\vect{v}_s$};
  \draw[Stealth-Stealth, thick] (2.55,0) -- (3.0,0);
  \node[above] at (2.78,0.05) {$\lambda'$};
  \draw[Stealth-Stealth, thick] (-3.0,0) -- (-1.95,0)
    node[midway, above] {$\lambda'$};
  \node[gray] at (0.15,0.55) {\scriptsize 过去位置};
\end{tikzpicture}

{\small 波源向右运动时的波前：每个圆的圆心都是它被
发出的位置（灰点），因此波峰在前方挤聚，$\lambda' = (v - v_s)T$，
在后方拉长，$\lambda' = (v + v_s)T$。}
\end{omfigure}

\begin{example}[用数字看救护车]\label{ex:g12:doppler-effect:siren}
一个警笛的频率为 $f = \qty{700}{Hz}$，以 $v_s = \qty{25}{m/s}$ 的速率驶近。
取 $v = \qty{340}{m/s}$，为空气中的传播速度（\cref{ch:g12:sound-acoustics}）。靠近时：
$f' = 700 \times 340/315 \approx \qty{756}{Hz}$；远离时：
$f' = 700 \times 340/365 \approx \qty{652}{Hz}$。经过时，音调下降了 \qty{104}{Hz}
——比值 $\approx 1.16$，约两个半音，任何耳朵都能分辨。
\end{example}

\begin{remark}[运动不改变什么]\label{rem:g12:doppler-effect:speed-unchanged}
波仍以 $v$ 的速度传播，只有介质决定
传播速度（\cref{ch:g12:mechanical-waves}）。波源的运动压缩的是\emph{波长}，而不压缩
速率。压缩也有极限：当 $v_s \to v$ 时，前方的波峰彼此堆叠，
$f' \to \infty$——超音速飞机所拖曳的激波，正是波前的一场“交通堵塞”。
\end{remark}

\section{运动的观察者与慢速近似}

\begin{proposition}[运动的观察者]\label{prop:g12:doppler-effect:moving-observer}
观察者以速率 $v_o$ 沿直线朝向（或远离）介质中静止的
波源运动。接收频率为
\[ f' = f\left(1 \pm \frac{v_o}{v}\right) \quad (+\ \text{靠近},\ -\ \text{远离}). \]
\end{proposition}

\begin{proof}
波形本身未受扰动：间距为 $\lambda = vT$ 的波峰以传播速度 $v$ 运动。
观察者迎着波峰奔跑，以相对速率 $v + v_o$ 遇到波峰，
每隔 $\lambda/(v + v_o)$ 秒遇到一个：$f' = (v + v_o)/\lambda = f(1 + v_o/v)$。背向奔跑时，相对速率为 $v - v_o$。
\end{proof}

\begin{definition}[径向速度]\label{def:g12:doppler-effect:radial}
波源相对于观察者的\emph{径向速度}\index{径向速度} $v_r$ 是
两者相对速度沿连线方向的分量——也就是
距离缩小或增大的快慢。只有这一分量会改变频率；
\emph{横穿}视线的运动在这一阶近似下不产生频率变化。
\end{definition}

\begin{proposition}[慢速运动：对任何人都适用的同一公式]\label{prop:g12:doppler-effect:slow}
当径向速率很小，$v_r \ll v$ 时，不再区分究竟是谁在运动：
\[
\frac{\Delta f}{f} = \frac{f' - f}{f} \approx +\frac{v_r}{v} \ \text{（靠近）},
\qquad
\frac{\Delta f}{f} \approx -\frac{v_r}{v} \ \text{（远离）}.
\]
\end{proposition}

\begin{proof}
观察者公式本身已恰为 $1 \pm v_o/v$。对于波源，记 $x = v_s/v$：
恒等式 $\frac{1}{1 - x} = 1 + x + \frac{x^2}{1 - x}$ 表明 $f'/f$ 与
$1 + x$ 只相差一个量级为 $x^2$ 的项——取 $x = 0.1$ 时，$1\%$ 的修正叠加在 $10\%$ 的偏移之上。
在一阶近似下，四个公式都归结为 $1 \pm v_r/v$。
\end{proof}

\begin{omfigure}
\begin{tikzpicture}
\begin{axis}[omaxis, width=8.8cm, height=4.8cm,
    xlabel={时间 $t$ (\unit{s})}, ylabel={频率 $f'$ (\unit{Hz})},
    xmin=-8, xmax=8, ymin=408, ymax=472,
    ytick={420,440,460}, xtick={-8,-4,0,4,8}]
  \addplot[dashed, omExo, domain=-8:8, samples=2] {460.3};
  \addplot[dashed, omExo, domain=-8:8, samples=2] {421.4};
  \addplot[dotted, omProp, thick, domain=-8:8, samples=2] {440};
  \addplot[omDef, thick, domain=-8:8, samples=120]
    {440*340/(340 + 15*(15*x)/sqrt((15*x)^2 + 64))};
\end{axis}
\end{tikzpicture}

{\small 一辆频率为 \qty{440}{Hz} 的警笛以 \qty{15}{m/s} 的速度驶过，距离麦克风 \qty{8}{m}
处时听到的频率平台值为 $fv/(v \mp v_s)$，而真实的 $f$（点线）
几乎在最近距离处经过，此时运动纯属横穿。}
\end{omfigure}

\begin{method}[读取经过数据]\label{met:g12:doppler-effect:pass}
一段运动波源经过时的录音显示两个平台：$f_1$（靠近）与 $f_2$
（远离）。则
\[
f = \frac{2 f_1 f_2}{f_1 + f_2}, \qquad
v_s = v\,\frac{f_1 - f_2}{f_1 + f_2}.
\]
事实上，$f_1 + f_2 = \frac{2fv^2}{v^2 - v_s^2}$，$f_1 - f_2 = \frac{2fv\,v_s}{v^2 - v_s^2}$：
两式相除得 $v_s$，再计算 $2f_1 f_2/(f_1 + f_2)$ 得 $f$。两个频率读数即可同时给出
真实的音调与速率——无需秒表，也无需卷尺。
\end{method}

=\section{用回声测量速率}

让波从运动目标上反射，多普勒效应便起两次作用。

\begin{proposition}[反射使频移加倍]\label{prop:g12:doppler-effect:double}
频率为 $f$ 的波射向以速率 $u \ll v$ 迎面靠近的反射体。
回波的频率为
\[ f' = f\,\frac{v + u}{v - u}, \qquad \text{于是} \qquad \Delta f \approx \frac{2u}{v}\,f . \]
\end{proposition}

\begin{proof}
两次频移串联发生。作为运动的\emph{观察者}，反射体接收到 $f(1 + u/v)$；
将所接收的波重新发出，此时成为运动的\emph{波源}，因而回波的频率为
$f' = f(1 + u/v)/(1 - u/v) \approx f(1 + 2u/v)$，当 $u \ll v$ 时。
\end{proof}

\begin{example}[雷达测速枪]\label{ex:g12:doppler-effect:radar}
一台交通雷达以 $f = \qty{24.15}{GHz}$ 的频率发射微波（$v = c = \qty{3.00e8}{m/s}$），
并将回波与发出的波叠加；拍（见\cref{ch:g12:sound-acoustics}）
直接发出拍频 $\Delta f = 2uf/c \approx \qty{161}{Hz}$（每 \unit{m/s} 汽车速率）
——从 \qty{24}{GHz} 的载波中分离出一个极易测量的音频频率。
因子 $2$ 不可省略：忘记它会使每位司机的车速被低估一半。
\end{example}

\begin{example}[多普勒超声]\label{ex:g12:doppler-effect:ultrasound}
一个医用探头向动脉发射频率 $f = \qty{5.0}{MHz}$ 的超声波
（在组织中 $v \approx \qty{1540}{m/s}$）；红细胞将其反射。血液以
$u = \qty{0.50}{m/s}$ 的速率流动时，回波的 $\Delta f = 2 \times 0.50 \times \num{5.0e6}/1540 \approx
\qty{3.2}{kHz}$ ——一个\emph{可听}的拍频：心脏科医生真能听到血液
在每次心跳时加速，然后测定 $u$，依据为 $\Delta f$。
\end{example}

\section{星光中的多普勒效应}

光是波（见\cref{ch:g12:light-as-wave}），其光谱中布满尖锐的谱线，
谱线的波长由每种化学元素确定（见\cref{ch:g10:light-spectra}）——一把
经实验室标定的标尺，印在每一颗恒星上。

\begin{definition}[红移与蓝移]\label{def:g12:doppler-effect:redshift}
若某波源的谱线出现的波长比在实验室中\emph{更长}，
则称该光谱发生\emph{红移}\index{红移}；若谱线的波长\emph{更短}，
则称其发生\emph{蓝移}\index{蓝移}——光谱向可见光谱的
红端或蓝端移动。
\end{definition}

\begin{proposition}[光的多普勒效应]\label{prop:g12:doppler-effect:light}
对以径向速率 $v_r \ll c$ 远离的波源，每个波长被观察到
伸长了
\[ \frac{\Delta \lambda}{\lambda} = \frac{\lambda' - \lambda}{\lambda} \approx \frac{v_r}{c} \quad \text{（红移）}, \]
并被同样程度地压缩，$\Delta\lambda/\lambda \approx -v_r/c$，对于靠近的
波源（蓝移）。
\end{proposition}
\admitted

\begin{remark}[准确公式的来源]\label{rem:g12:doppler-effect:honest}
光不需要介质，因此不能照搬声波的推导——不存在“在介质中静止”
的说法。精确公式来自狭义相对论，本学年稍后将学到
（见\cref{ch:g12:special-relativity}）；修正项的量级为 $(v_r/c)^2$，对飞机、恒星和近邻星系
而言可以忽略。以下放心使用慢速近似公式。
\end{remark}

\begin{omfigure}
\begin{tikzpicture}[font=\small]
  \foreach \y/\xl/\c/\t in {2.4/1.8/black/实验室，
      1.2/0.6/omDef/靠近的恒星，0/3.72/omThm/远离的恒星} {
    \fill[gray!18] (0,\y) rectangle (6,\y+0.55);
    \draw[\c, very thick] (\xl,\y) -- (\xl,\y+0.55);
    \node[right] at (6.15,\y+0.28) {\t};
  }
  \draw[dashed, gray] (1.8,-0.15) -- (1.8,2.95);
  \draw[-Stealth] (-0.2,-0.55) -- (6.4,-0.55)
    node[below left] {波长 $\lambda$ (\unit{nm})};
  \foreach \x/\l in {0/656.0, 3/656.5, 6/657.0}
    \draw (\x,-0.48) -- (\x,-0.62) node[below] {\l};
\end{tikzpicture}

{\small 同一条氢谱线在三种光谱中的情况：实验室中为 \qty{656.30}{nm}，
靠近的恒星发生蓝移，远离的恒星发生红移。对照虚线参考线读出位移，
即可求得径向速度。}
\end{omfigure}

\begin{example}[测算恒星的逃逸]\label{ex:g12:doppler-effect:vega}
在明亮恒星织女星的光谱中，实验室波长为
\qty{656.30}{nm} 的氢谱线测得波长为 \qty{656.27}{nm}：$\Delta\lambda = \qty{-0.03}{nm}$，发生蓝移，
因此织女星正在靠近，其速率为 $v_r = c\,|\Delta\lambda|/\lambda \approx \qty{1.4e4}{m/s}$
——即 \qty{14}{km/s}，由约两万分之一的位移读出。
\end{example}

\begin{example}[摇摆的恒星：双星与系外行星]\label{ex:g12:doppler-effect:wobble}
当两颗恒星相互绕转时，它们的谱线会周期性地交替蓝移与红移——这就是
\emph{分光双星}\index{分光双星}，无需分辨这对恒星，即可测得其轨道。行星对恒星的作用也是如此，只是微弱得多：二者都绕其共同质心运动，
因此恒星的径向速度以几十 \unit{m/s} 的幅度振荡——对应的位移为
$\Delta\lambda/\lambda \sim 10^{-7}$。
正是通过探测这一位移，人们发现了第一颗类太阳恒星周围的行星；周末问题将重现这一发现，
包括其中的数字。
\end{example}

\begin{remark}[星系红移]\label{rem:g12:doppler-effect:expansion}
遥远星系的光谱全部发生\emph{红}移，且越远的星系红移越明显：宇宙中的距离正在拉伸。对于近邻星系，
$v_r = c\,\Delta\lambda/\lambda$ 可读出退行速率；对于最遥远的星系，
位移已超出慢速公式可以可靠处理的范围，
相关的推算将从
\cref{ch:g12:special-relativity}起继续，并在大学分册中展开。
\end{remark}

\section{练习}

\begin{exercise}[$\star$]\label{exo:g12:doppler-effect:1}
一辆消防车的警笛发出 \qty{700}{Hz} 的声音，消防车以 \qty{25}{m/s} 的速度行驶（$v = \qty{340}{m/s}$）。
求行人听到的频率：(a) 位于车前方时；(b) 位于车后方时。
(c) 司机听到的频率是多少？
\end{exercise}

\begin{exercise}[$\star$]\label{exo:g12:doppler-effect:2}
一位朋友说：“救护车离开时音调降低，是因为声音随距离增大而变弱。”
请厘清这里混淆的两种效应，并说明二者各自取决于什么。
\end{exercise}

\begin{exercise}[$\star$]\label{exo:g12:doppler-effect:3}
一名骑车人以 \qty{6.0}{m/s} 的速度径直朝静止的警笛行驶，该警笛发出 \qty{700}{Hz} 的声音。
他听到的频率是多少？若他径直远离警笛，频率又是多少？
应使用哪一个公式？为何不能使用运动声源的公式？
\end{exercise}

\begin{exercise}[$\star$]\label{exo:g12:doppler-effect:4}
一列火车的汽笛发出 \qty{400}{Hz} 的声音，火车以 \qty{30}{m/s} 的速度行驶。求火车前方、后方及静止时的波长，
并指出哪位观察者接收到哪一种波长。
\end{exercise}

\begin{exercise}[$\star$]\label{exo:g12:doppler-effect:5}
一部手机播放 \qty{440}{Hz} 的音调，由一名以 \qty{5.0}{m/s} 速度奔跑的人携带。用慢速公式估算迎面走来的人所听到的频移。
将其与半音（频率约为 $6\%$）进行比较：
音乐家会察觉到吗？
\end{exercise}

\begin{exercise}[$\star\star$]\label{exo:g12:doppler-effect:6}
一只蝙蝠以 \qty{6.0}{m/s} 的速度径直飞向墙壁，发出 \qty{50.0}{kHz} 的超声波。证明它听到的回声频率为 $f' = f(v + v_b)/(v - v_b)$，
再求 $f'$ 及其频移。
为什么这实际上就是雷达测速枪的公式？
\end{exercise}

\begin{exercise}[$\star\star$]\label{exo:g12:doppler-effect:7}
一台工作频率为 \qty{24.15}{GHz} 的雷达测速枪，测得靠近的汽车产生的拍频为 \qty{4.0}{kHz}。
求汽车的速率，分别以 \unit{m/s} 和 \unit{km/h} 表示。
若一辆汽车恰以 \qty{50}{km/h} 行驶，它会产生怎样的拍频？
\end{exercise}

\begin{exercise}[$\star\star$]\label{exo:g12:doppler-effect:8}
一台沿动脉方向瞄准的 \qty{5.0}{MHz} 多普勒超声探头（在组织中 $v = \qty{1540}{m/s}$）
测得拍频为 \qty{2.6}{kHz}。求血流速率。为什么探头必须
\emph{沿着}血流方向放置，而不是与之垂直？
\end{exercise}

\begin{exercise}[$\star\star$]\label{exo:g12:doppler-effect:9}
取 $f = \qty{1000}{Hz}$，$v = \qty{340}{m/s}$，接近速率为 \qty{34}{m/s}。
分别精确计算：(a) 声源运动时，(b) 观察者运动时的 $f'$。
将两者与慢速预测 $f(1 + v_r/v)$ 比较，并解释二者差异的大小。
\end{exercise}

\begin{exercise}[$\star\star$]\label{exo:g12:doppler-effect:10}
记录一辆滑板车经过麦克风时的图像，显示出 \qty{465}{Hz} 与 \qty{415}{Hz} 两个平台。
利用\cref{met:g12:doppler-effect:pass}，求喇叭的真实频率，以及滑板车以 \unit{km/h} 表示的速率。
为什么真实频率\emph{不是}这两个数值的平均值？
\end{exercise}

\begin{exercise}[$\star\star$]\label{exo:g12:doppler-effect:11}
在某恒星的光谱中，实验室波长为 \qty{656.3}{nm} 的氢谱线被观测到位于 \qty{656.5}{nm} 处。
是红移还是蓝移？求该恒星的径向速度及其方向。
\end{exercise}

\begin{exercise}[$\star\star\star$]\label{exo:g12:doppler-effect:12}
一颗恒星因一颗看不见的行星而以 \qty{55}{m/s} 的速度摆动。求波长为 \qty{500}{nm} 的谱线，其波长摆动的幅度，
分别以长度和占 $\lambda$ 的比例表示。
数量级分析：为什么行星探测
必须等到光谱仪稳定度达到 $10^{7}$ 分之一才能实现？
\end{exercise}

\begin{exercise}[$\star\star\star$]\label{exo:g12:doppler-effect:13}
一个星系光谱中的每条谱线都被拉伸了 $2.4\%$。求其退行速率。
天文学家发现，这样的速率在天空各个方向上都与距离成正比：这暗示了怎样的图景？
慢速公式在此处是否仍然可靠？
\end{exercise}

\begin{exercise}[$\star\star\star$]\label{exo:g12:doppler-effect:14}
太阳在自转：其圆面的一边以约 \qty{2.0}{km/s} 的速度靠近我们，另一边则以同样的速度远离。
这对来自整个日面的 \qty{656.3}{nm} 谱线有何影响——是使其位移，还是使其展宽？
求该效应以 \unit{nm} 表示的大小。
\end{exercise}

\begin{exercise}[$\star\star\star$]\label{exo:g12:doppler-effect:15}
一辆警车以 \qty{40}{m/s} 的速度行驶，警笛频率为 \qty{700}{Hz}，
正在追赶一辆同向行驶、速率为 \qty{30}{m/s} 的卡车。(a) 试从卡车司机的角度论证 $f' = f(v - v_o)/(v - v_s)$。
(b) 计算 $f'$。(c) 证明：若卡车的速率与警车相同，频移将消失，
并说明这为何是理所当然的。
\end{exercise}

\section{问题：寻找系外行星}

\begin{problem}[{周末问题——寻找系外行星：火车鸣笛校准听觉，雷达枪预演双重频移，
到周日夜晚，一颗恒星光谱的摆动仅为五百万分之一，
却称出了一颗无人见过的行星}]
\label{pb:g12:doppler-effect:1}
一个天文社团用一个周末研究一个想法——多普勒频移——，先在地球上预演，
再瞄准一颗类太阳恒星。数据：声音在空气中的速率 $v = \qty{340}{m/s}$，
$c = \qty{3.00e8}{m/s}$，$G = \qty{6.67e-11}{N.m^2/kg^2}$；恒星质量为
$M = \qty{2.0e30}{kg}$。在本阶段不加证明地承认，由圆轨道的力学
（\cref{ch:g12:satellites-kepler}）可知，质量为 $m$ 的行星在半径为 $r$ 的圆轨道上运动，
其周期与速率由 $r^3 = \frac{G M P^2}{4\pi^2}$ 和 $V = \frac{2\pi r}{P}$ 给出，
而恒星绕共同质心运动，其摆动速率为 $K = \frac{m}{M}\,V$。

\medskip
\noindent\textbf{第一部分——星期五：平交道口。}
一列火车以恒定速率驶过，鸣着笛；一部手机在火车靠近时记录到 $f_1 = \qty{471}{Hz}$，
驶过之后记录到 $f_2 = \qty{411}{Hz}$。
\begin{enumerate}
  \item 借助波前密集的图像解释：为何 $f_1 > f_2$，尽管鸣笛的频率从未改变。
  \item 写出联系 $f_1$、$f_2$ 与真实频率 $f$ 和火车速率 $u$ 的两个方程。
  \item 证明 $u = v\,\dfrac{f_1 - f_2}{f_1 + f_2}$，并以 \unit{m/s} 和 \unit{km/h} 计算其值。
  \item 证明 $f = \dfrac{2 f_1 f_2}{f_1 + f_2}$，并计算其值。为何 $f$ 略
        \emph{低于} $f_1$ 与 $f_2$ 的算术平均值？
  \item 计算火车前方与后方的波长，以及静止波长。
\end{enumerate}

\medskip
\noindent\textbf{第二部分——星期六：雷达枪。}
一支交通巡逻队演示其雷达：$f = \qty{24.125}{GHz}$，雷达枪读取
出射波与回波之间的拍频。
\begin{enumerate}[resume]
  \item 解释为何运动汽车的回波被\emph{两次}频移——指出汽车在每次频移中
        各扮演的角色。
  \item 证明：当汽车速率 $u \ll c$ 时，拍频为 $\Delta f \approx \dfrac{2u}{c}\,f$。
  \item 某汽车的回波对应 $\Delta f = \qty{3.55}{kHz}$，求 $u$。
  \item 换算为 \unit{km/h}，并与当地标示的 \qty{80}{km/h} 限速比较。
  \item \qty{1}{km/h} 对应多少赫兹的拍频？请讨论：
        是什么使 \qty{24}{GHz} 波的这种微小相对频移易于测量？
\end{enumerate}

\medskip
\noindent\textbf{第三部分——星期日：摆动的恒星。}
社团下载了一颗类太阳恒星逐夜测得的径向速度，
图示如下。

\begin{omfigure}
\begin{tikzpicture}
\begin{axis}[omaxis, width=8.8cm, height=4.6cm,
    xlabel={$t$（天）}, ylabel={$v_r$（\unit{m/s}）},
    xmin=0, xmax=8.4, ymin=-72, ymax=72,
    xtick={0,2.1,4.2,6.3,8.4},
    xticklabels={$0$,$2.1$,$4.2$,$6.3$,$8.4$}, ytick={-55,0,55}]
  \addplot[dashed, omExo, domain=0:8.4, samples=2] {55};
  \addplot[dashed, omExo, domain=0:8.4, samples=2] {-55};
  \addplot[omDef, thick, domain=0:8.4, samples=140] {55*sin(360*x/4.2)};
\end{axis}
\end{tikzpicture}

{\small 恒星十二个夜晚的径向速度：一条干净的正弦曲线——
是圆形轨道上、侧视所见天体的特征印记。}
\end{omfigure}

\begin{enumerate}[resume]
  \item 逐夜在恒星的\emph{光谱}中实际测量的是什么？
        写出用来将其换算为 $v_r$ 的关系式。
  \item 由曲线读出摆动的振幅 $K$ 与周期 $P$。
  \item 计算一条 \qty{550}{nm} 谱线对应的摆动量 $\Delta\lambda$，以及
        $\Delta\lambda/\lambda$。与第一部分比较：此测量
        比火车的测量难多少？
  \item 看不见的行星为何会使恒星摆动？正弦形状
        说明轨道的什么性质？
  \item 为何此方法只测量恒星运动的\emph{径向}分量，以及
        若轨道面正对我们，将看到什么？
\end{enumerate}

\medskip
\noindent\textbf{第四部分——星期日夜晚：称量不可见之物。}
设轨道为圆形且侧视，于是 $K$ 即恒星的全部轨道速率。
\begin{enumerate}[resume]
  \item 由 $P$ 及已承认的关系式 $r^3 = G M P^2/(4\pi^2)$，求行星的
        轨道半径 $r$。将其与日地距离 \qty{1.5e11}{m} 比较。
  \item 计算行星的轨道速率 $V = 2\pi r/P$。
  \item 由 $K = (m/M)\,V$，求行星的质量 $m$。
  \item 将 $m$ 与木星（$\qty{1.9e27}{kg}$）及地球（$\qty{6.0e24}{kg}$）的质量相比，并
        将 $r$ 与日水距离（$\qty{5.8e10}{m}$）相比。这是一个怎样的世界？
  \item 若轨道倾斜而非侧视，真实质量会比你算出的值
        更大还是更小？用一句话总结：一个仅为五百万分之一的周期性频移，
        究竟给出了什么？
\end{enumerate}
\end{problem}
